\providecommand{\E}{\mathbb E}
\subsection{Model, inputs, and arithmetic scope}
Write $\theta=(U,W)$, with all input rows followed by all output rows,
$f_\theta(x)=\sum_iw_i(u_i^Tx)_+$ and
$\mathcal L=\frac12\mathbb E_n\|x-f_\theta(x)\|^2$.
The full training residual is $e=x-f_\theta(x)$, and $F=J^*e$ denotes a
compatible gradient field. The adjoint uses the empirical inner product.
The bounds below cover all compatible selectors; no global uniqueness is
assumed. Global Clarke flows and the compatible representation are those
of the supplied revised appendix's nonsmooth-flow lemma.

Put $r=\sin(\pi/36)$, $k=\cos(\pi/36)$, $v=(r,-1)$, and
$M=\sum_{i\in I}w_iu_i^T$. For $Y=e_1^TMv-r$, the empirical quantity is
\[
 Q_G=Y(r\dot m_{11}-\dot m_{12})+r\dot m_{21}-\dot m_{22}
     =\nabla\Phi_I(\theta)^TF(\theta),\qquad
 \Phi_I=\tfrac12Y^2+e_2^TMv.
\]
Its output gradient is $\ell=(Y,1)$, so
$g_{u_i}=\mathbf1_Iv(w_i^T\ell)$ and
$g_{w_i}=\mathbf1_I\ell(u_i^Tv)$.
The unit central probe is $(r,-k)$; no $k$ factors are inserted into $Q_G$.
The canonical script uses exact field evaluation. This is distinct from the
finite-difference estimator used by the nine-setting study.

The proof inputs and code are hashed in the accompanying manifest. The
repository commit is
\begin{center}\texttt{ae16f9c0e4e51269c624ffb196184568f2cfe6d9}.\end{center}
The exact archived initial parameter rows are preserved; freshly evaluating
the trigonometric initialization changes some entries and is not used.
The supplied binary dataset exactly matches the committed seeded generator
under the recorded runtime. The historical state archive contains neither
raw $X$ nor an independent $X$ hash, so historical bitwise dataset identity
is not independently established. The theorem is about the supplied binary
instance, not an iid Gaussian event.

Arithmetic bounds use midpoint-radius intervals, outward IEEE rounding
bounds for sums and dot products, and interval transcendental evaluation.
They assume standard round-to-nearest operations without unsafe fast-math.
All numerical tube radii are verified by strict inequalities; no integrator
accuracy claim is accepted without a reference-defect bound.

\subsection{A larger certified entry region}
Let $X$ be the exact supplied binary dataset with $N=4000$, $h=200$,
and the canonical generation parameters $(\mu_1,\mu_2,\sigma)=(3,2,0.15)$.
Let $\widehat\theta_{3.4}$ be the supplied reference state. The fixed
$71$-neuron cohort $I$ is selected by the archived in-distribution rule at
$t=4$, not by OOD outcomes. Put
\[
 \xi_\alpha=(\cos\alpha,\sin\alpha),\qquad
 \mathcal I=\{\alpha:|\alpha+17\pi/36|<10^{-6}\}.
\]
The theorem refers to global Clarke gradient flows of the empirical loss;
it applies to every such flow satisfying its entry condition. No global
uniqueness assumption is needed. Compatible selector fields are used as an
overapproximation of the Clarke inclusion, including at zero preactivation.

\begin{theorem}[Canonical continuation from a learned-state neighborhood]
\label{initcert:wide-theorem}
If
\begin{equation}\label{initcert:entry}
 \|\theta(3.4)-\widehat\theta_{3.4}\|_2\le1.2\cdot10^{-5},
\end{equation}
then, uniformly for $\alpha\in\mathcal I$,
\begin{align*}
 E(\xi_\alpha,3.4)-E(\xi_\alpha,3.7)&>1.2\cdot10^{-4},\\
 E(\xi_\alpha,3.9)-E(\xi_\alpha,3.7)&>4\cdot10^{-5}.
\end{align*}
The error has an interior minimum on $[3.4,3.9]$. At entry the cohort has
$m_{11}>0.9323$, input angles below $15^\circ$, and output angles below
$20^\circ$. For the central probe, both $Q_G$ and the full OOD rate are
negative at entry and positive at $3.9$ for all compatible training selectors.
The strong-cluster contribution to the full OOD rate is negative at both
endpoints, the weak-cluster contribution is positive, and the latter
dominates at the late endpoint.
Moreover, almost everywhere, for all probes in the sector,
\[
 \dot E<-0.0012\quad\hbox{on }[3.4,3.40001],\qquad
 \dot E>0.00040\quad\hbox{on }[3.9,3.90001].
\]
The sector has width $2\cdot10^{-6}$ radians. Neither an initialized
entry guarantee nor a unique continuous zero of $Q_G$ is asserted.
\end{theorem}

\begin{proof}
The certificate applies the one-sided tube inequality and the verified-step
criterion below to $2500$ reference intervals of exact length $1/5000$.
Every reference coefficient is an exact stored binary number; its numerical
integration provenance is not a proof premise. The uniform interval bounds
are independent of the chosen initial error within the proposed tubes.
An independent scalar replay with upward-rounded initial radius
$1.2\cdot10^{-5}$ verifies every strict tube inequality. At the snapshots,
\[
 \|\theta(3.7)-\widehat\theta_{3.7}\|<2.828\cdot10^{-5},\qquad
 \|\theta(3.9)-\widehat\theta_{3.9}\|<4.251\cdot10^{-5}.
\]
The full step ledger is \texttt{WIDE\_LOCAL\_FLOW\_ENCLOSURE.csv}.
Static bounds use radii $2\cdot10^{-5}$, $3.5\cdot10^{-5}$ and
$4.6\cdot10^{-5}$ at the three snapshots. On their respective balls,
\begin{align*}
 E(\xi_0,3.4)&\in[0.48314641,0.48315250],\\
 E(\xi_0,3.7)&\in[0.48298853,0.48299930],\\
 E(\xi_0,3.9)&\in[0.48305831,0.48307257],
 \qquad \xi_0=\xi_{-17\pi/36}.
\end{align*}
These displayed intervals are enlarged from the computed enclosures.
The full-precision central gaps exceed $0.00014712$ and $0.00005902$.
On the unit circle the error is Lipschitz in the probe with constant
$(1+\|\theta\|^2/2)^2<9.084$. Transporting two errors through an angular
change below $10^{-6}$ costs less than $1.8168\cdot10^{-5}$.
The resulting verified sector gaps are greater than $0.00012895$ and
$0.00004085$, proving the stated margins and the interior minimum.

The endpoint interval bounds for $10^3$ times the exact rates are:
\begin{center}
\begin{tabular}{llrr}
\toprule
Statistic & contribution & $t=3.4$ & $t=3.9$\\
\midrule
$Q_G$ & total & $[-1.140,-0.960]$ & $[0.630,1.029]$\\
$\dot E$ & total & $[-1.379,-1.215]$ & $[0.406,0.771]$\\
$\dot E$ & strong cluster & $[-2.468,-2.320]$ & $[-1.273,-0.939]$\\
$\dot E$ & weak cluster & $[1.055,1.138]$ & $[1.601,1.788]$\\
\bottomrule
\end{tabular}
\end{center}
Each cluster uses the same full-network residual and normalization $1/N$;
this is a signed field decomposition, not a retraining ablation.
The same interval calculation verifies the mass and cone statements.

The speed bounds on the first and last static balls are $0.385$ and
$0.309$. The verified endpoint radii leave enough room for a forward
extension of $10^{-5}$ at each endpoint. A first-exit argument keeps the
flow in those balls. The probe gates remain constant across each static
ball and sector. Indeed their central score margins exceed
$3.7575\cdot10^{-5}$ and $5.6601\cdot10^{-5}$, respectively;
subtracting the parameter radius and an additional $2.01\cdot10^{-6}$
for sector transport still leaves positive margins.
With $P<2.01$, the OOD rate varies with the probe by at most
$2(1+P^2/2)P\|F\|\|\xi-\xi_0\|<5\cdot10^{-6}$.
Subtracting this cost from the full-precision rate margins gives the two
uniform rate intervals. The exact static reports and arithmetic contract
are supplied with the certificate.
\end{proof}

\paragraph{Scope.}
Condition~\eqref{initcert:entry} specifies a parameter neighborhood of an
ID-trained state; it does not assume an OOD sign. Its radius is twelve times
the earlier sufficient radius. Membership of the canonical initialized flow
in this neighborhood remains unproved. The result therefore does not
strengthen the paper's existing random-initialization theorem or establish
that specialization alone causes reversal. The offline cohort and witness
times are retrospective, and the sector is deliberately conservative.

\subsection{Static sign certificates with all training-gate changes}
\label{initcert:static}
This section concerns parameter neighborhoods, not their reachability from
initialization. Its dataset consists of the exact binary numbers in the
accompanying \texttt{canonical\_witness\_states.npz}. Rounding a Gaussian
sample is not an iid Gaussian theorem.

Let $y$ be a parameter center, $R>0$, $P=\|y\|$, $\bar P=P+R$, and
$D=PR+R^2/2$. A cluster index set $C$ uses the full-dataset normalization
$\|z\|_C^2=N^{-1}\sum_{n\in C}\|z_n\|^2$. Put
$\lambda_C\ge\|N^{-1}\sum_{n\in C}x_nx_n^T\|$.
Write $F_C$ for that cluster's contribution to the full gradient field,
using the \emph{full-network residual}; $F=F_1+F_2$.
Let $F_C^0$ denote its polynomial extension with the training gates fixed
at $y$. No fixed-gate hypothesis is imposed on the actual field in the ball.

For either potential used here, set $g=\nabla\Phi$, $G_0=\|g(y)\|$.
For the empirical statistic use the potential $\Phi_I$ defined above. For the full derivative, fix the actual
probe-active set at $y$ and use
$\Phi(\theta)=\frac12\|M_{\rm act}(\theta)\xi-\xi\|^2$.
This second choice is valid throughout the ball if
$\min_i|u_i(y)^T\xi|>R\|\xi\|$.
In either case let $v$ be its input vector, $I$ its fixed index set,
$P_I=\|(u_i,w_i)_{i\in I}(y)\|$, and $\ell_y$ its output gradient.
Define
\begin{align*}
 L_C&=\lambda_C\bar P^2+
 \sqrt{\lambda_C}\big(\|e(y)\|_C+\sqrt{\lambda_C}D\big),\\
 H_g&=\|v\|^2(P_I+R)^2+
 \|v\|\big(\|\ell_y\|+\|v\|(P_IR+R^2/2)\big).
\end{align*}

Here are explicit bounds for the training-gate correction. Put
$q_{ni}=u_i(y)^Tx_n$, $s_n=\|x_n\|$, $e_n=e(y;x_n)$, and
\begin{align*}
 d_{ni}&=(Rs_n-|q_{ni}|)_+,&
 p_{ni}&=\mathbf1_{|q_{ni}|\le Rs_n},\\
 b_{ni}&=|w_i(y)^Te_n|+R\|e_n\|
                 +(\|w_i(y)\|+R)Ds_n,\\
 j_i^u&=N^{-1}\sum_{n\in C}p_{ni}b_{ni}s_n,&
 j_i^w&=N^{-1}\sum_{n\in C}d_{ni}(\|e_n\|+Ds_n),\\
 z_n&=\sum_i(\|w_i(y)\|+R)d_{ni},&
 B_C&=\Big[\sum_i((j_i^u)^2+(j_i^w)^2)\Big]^{1/2}
                  +\sqrt{\lambda_C}\bar P\|z\|_C.
\end{align*}
Every neuron is retained in these quantities.

\begin{proposition}[A static neighborhood enclosure]\label{initcert:static-bound}
For any compatible training selectors and $\|\theta-y\|\le R$,
\begin{equation}\label{initcert:static-equation}
 \left|g(\theta)^TF_C(\theta)-g(y)^TF_C(y)\right|
 \le R\left[H_g(\|F_C(y)\|+L_CR)+G_0L_C\right]
                   +(G_0+H_gR)B_C.
\end{equation}
This applies to the total derivative and separately to each cluster's
signed contribution. Cluster contributions always use the same full residual.
\end{proposition}
\begin{proof}
The Jacobian norm on cluster $C$ is at most
$\sqrt{\lambda_C}\|\theta\|$: apply Cauchy--Schwarz across neuron and
input/output parameter coordinates before taking the empirical average.
Integration along a parameter segment gives
$\|f_\theta-f_y\|_C\le\sqrt{\lambda_C}D$ and the pointwise bound
$\|f_\theta(x_n)-f_y(x_n)\|\le Ds_n$.
The same bounds hold for the fixed-mask polynomial extension.
Its field derivative is $-J_C^*J_C+H_{e,C}$, where the neuronwise
residual Hessian has norm at most
$\sqrt{\lambda_C}\|e\|_C$. Thus
$\|DF_C^0\|\le L_C$. Differentiating either potential gives
$\|Dg\|\le H_g$ on the ball. Adding and subtracting
$g(y)^TF_C^0(\theta)$ proves the first term of \eqref{initcert:static-equation}.

A changed training activation differs from its fixed-mask extension by at
most $d_{ni}$, and its selector can differ only when $p_{ni}=1$.
The pointwise residual and weight bounds give
$|w_i(\theta)^Te_\theta(x_n)|\le b_{ni}$. The direct input- and
output-gradient changes are therefore bounded by $j_i^u,j_i^w$.
The change of residual between the actual network and the fixed-mask
extension has norm at most $z_n$. Applying the fixed-mask adjoint to
this residual costs at most $\sqrt{\lambda_C}\bar P\|z\|_C$.
Consequently $\|F_C-F_C^0\|\le B_C$. Multiply by
$\|g(\theta)\|\le G_0+H_gR$ to finish. Endpoint partial selectors obey
the same bounds.
\end{proof}


% Research module: real-arithmetic bounds. Floating-point validation is separate.
\subsection{An enclosure through training-gate changes}
This section states the real-arithmetic stability inequality used by the
interval validator. A successful floating-point run is not by
itself a certificate; its input and rounding enclosures must also be verified.

Write $F(\theta)=J_\theta^*(x-f_\theta)$ for the selected vector field.
Fix a reference state $y$ and a radius $R$. Assume that along every straight
parameter segment within this ball, the smooth pieces obey
\[
 \max_i\|\E_n[e_\theta(x)x^T\alpha_i(x)]\|\le a.
\]
For each neuron and sample choose an upper bound
$h_{in}\ge N^{-1}\|x_n\|(w_i^Te_\theta(x_n))_+$ throughout the ball.
Set $d_{in}=|u_i(y)^Tx_n|/\|x_n\|$ for nonzero samples, and
\[
 B_i(s)=\sum_{n:d_{in}\le s}h_{in},\qquad 0\le s\le R.
\]
If the center ranges along a reference subinterval, lower bounds for the
margins and uniform upper bounds for the weights can replace these quantities.

\begin{lemma}[One-sided stability with adverse gate jumps]\label{initcert:tube}
Suppose $B_i(s)\le a_i+\beta s$ for $0\le s\le R$, where $a_i,\beta\ge0$.
Then for $\|z\|\le R$,
\[
 \langle z,F(y+z)-F(y)\rangle
 \le(a+\beta)\|z\|^2+\|(a_i)_i\|_2\|z\|.
\]
Compatible endpoint selectors are allowed. Along an absolutely continuous
reference path $y(t)$ with defect $\|F(y(t))-\dot y(t)\|\le\eta$, the
error $r(t)=\|\theta(t)-y(t)\|$ therefore satisfies, while $r\le R$,
\[
 D^+r\le(a+\beta)r+\|(a_i)_i\|_2+\eta
\]
in the comparison sense (or almost everywhere away from $r=0$).
\end{lemma}
\begin{proof}
On a smooth activation region,
$DF=-J^*J+H_e$. The residual Hessian $H_e$ is block diagonal across
neurons, with each block having off-diagonal entries $A_i,A_i^T$, where
$A_i=\E_n[e x^T\alpha_i]$. Thus $\lambda_{\max}(DF)\le a$.
Integrate this inequality along the segment $y+sz$, separating its finitely
many training hyperplane crossings. When one gate crosses, the vector-field
jump has only an input-coordinate component
$\Delta F_{u_i}=N^{-1}\Delta\alpha_i(w_i^Te_x)x$; the output-coordinate
component is continuous because the preactivation is zero there. Its pairing
with $z$ is at most
$N^{-1}(w_i^Te_x)_+|z_{u_i}^Tx|\le h_{in}\|z_{u_i}\|$.
Only samples with $d_{in}\le\|z_{u_i}\|$ can cross. The sum of jumps is
therefore bounded by
\[
 \sum_i\|z_{u_i}\|B_i(\|z_{u_i}\|)
 \le\|(a_i)_i\|_2\|z\|+\beta\|z\|^2.
\]
Simultaneous crossings add; endpoint partial selectors satisfy the same
upper bound by taking partial jumps. Segments that lie in a hyperplane
contribute no conormal pairing for that gate. This proves the first claim.
Take the inner product of the error equation with $\theta-y$ and divide
by $r$ when it is positive. Standard norm regularization gives the comparison
at zero.
\end{proof}

For a finite set of distances, permissible intercepts are explicit:
\[
 a_i(\beta)=\max\big\{0,\max_{0\le s\le R}(B_i(s)-\beta s)\big\}.
\]
The inner maximum only needs the sorted sample distances, including zero.
There is no density assumption or asymptotic substitution for these counts.

\begin{corollary}[A verified step criterion]\label{initcert:tube-step}
On a reference interval of length $h$, let the preceding bounds be uniform
on its radius-$R$ tube. Put $k=a+\beta$ and
$\omega=\|(a_i)_i\|_2+\eta$. If the initial error is at most $r_0<R$ and
\[
 e^{kh}r_0+\omega\frac{e^{kh}-1}{k}<R,
\]
then the flow remains in the tube on the whole interval and its terminal
error is bounded by the left side. At $k=0$ the quotient is interpreted as
$h$. Successive verified steps give an enclosure from the initial state.
\end{corollary}
\begin{proof}
Apply the scalar integrating-factor inequality until the first exit. Its
upper bound is nondecreasing in time and strictly smaller than $R$ at the
terminal time, excluding a first exit. Iterate using the terminal upper bound
as the next initial radius.
\end{proof}
The absolute reference defect must also include any gates crossed by the
reference path. The favorable sign used in Lemma~\ref{initcert:tube} applies to
\emph{stability}, not automatically to the defect of a numerical interpolant.
Using only a smooth-region defect bound would omit this term and would not
constitute a certificate.

\subsubsection{Cubic-reference defect used in the local validator}
Let $y(s)$ be the cubic Hermite interpolant of two stored parameter states
and two stored velocities, over a step of exact length $1/5000$.
The stored velocities are merely coefficients of the reference polynomial;
they need not be exact field values. Put $a=1/10000$ and center the step
at $s=0$. Fix the actual training masks at $y(0)$ and denote that polynomial
field by $F^0$. The validator encloses
\begin{align*}
 d_0&=\|F^0(y)-y'\|,\\
 d_1&=\|DF^0(y)y'-y''\|,\\
 d_2&=\|D^2F^0(y)[y',y']+DF^0(y)y''-y'''\|
\end{align*}
at the midpoint. Suppose $P_*,V_*,A_*,T_*$ bound respectively
$\|y\|,\|y'\|,\|y''\|,\|y'''\|$ throughout the step, and $L_*$
bounds $\|DF^0\|$. For the covariance bound $\lambda$,
\[
 \|D^2F^0\|\le3\lambda P_*,\qquad
 \|D^3F^0\|\le3\lambda.
\]
Indeed the fixed-mask output is quadratic with second derivative norm
at most $\sqrt\lambda$, while $\|J\|\le\sqrt\lambda P_*$;
differentiate $F^0=J^*e$ twice and three times.
Taylor's integral remainder therefore gives
\[
 \eta_{\rm sm}=d_0+ad_1+\tfrac12a^2d_2+
 \tfrac16a^3\big(3\lambda V_*^3+9\lambda P_*V_*A_*+L_*T_*\big)
\]
as an upper bound for the fixed-mask defect on the whole step.
The third derivative of the defect uses $y''''=0$.

The true defect additionally includes $\eta_{\rm gate}$, the difference
between the actual field and $F^0$ along the reference polynomial.
It is bounded by the direct field argument in
Proposition~\ref{initcert:static-bound}, using neuronwise input and output excursions
\[
 d_i^u=a\|y'_{u_i}(0)\|+\tfrac12a^2\|y''_{u_i}(0)\|
                    +\tfrac16a^3\|y'''_{u_i}\|,
\]
and the analogous $d_i^w$. The corresponding whole-parameter excursion
bounds residual changes. Thus $\eta=\eta_{\rm sm}+\eta_{\rm gate}$ is
uniform on the step, including all reference-path gate changes.

For the one-sided rate on a tube, let $A_i^0$ be its midpoint residual
matrices and $D_*$ the pointwise output-change coefficient from the midpoint
throughout the tube. If $p_{in}$ flags every potentially changed gate there,
a valid smooth-piece bound is
\[
 a_* = \max_i\|A_i^0\|+\lambda D_*
       +\max_i N^{-1}\sum_n p_{in}\|e_y(x_n)\|\|x_n\|.
\]
For gate jumps the validator bounds the positive conormal residual throughout
the same tube, sorts lower bounds on its reference-path gate distances,
and sums upward-rounded weights. With $\beta=1$, the intercept
$\max_s(B_i(s)-s)_+$ only requires those finite sorted distances.
All tube radii are explicit inputs, rather than inferred from successful
numerical integration. Corollary~\ref{initcert:tube-step} is checked at every
step. Its strict inequality is what validates the reference, independently
of the integration method originally used to choose its coefficients.
